\section{Trial boot, health, and confirmation}

The trial image starts a confirmation controller as early as practical. It
reports \code{TRIAL\_STARTED} only after the network stack is available, but
network availability alone is not a sufficient health signal.

The minimum health gate includes:

\begin{itemize}
  \item verified boot and required filesystem mounts;
  \item successful schema/data migration or verified backward-compatible state;
  \item watchdog stability for the configured observation interval;
  \item startup of all mandatory services and device drivers;
  \item access to critical sensors, actuators, and secure storage; and
  \item no platform-defined fatal diagnostic or safety condition.
\end{itemize}

Cloud reachability \should be a confirmation requirement only when the device can
reliably distinguish a firmware defect from an external outage. Otherwise the
device confirms locally and reports when connectivity returns.

Confirmation is an atomic boot-control operation. The device \must first ensure
the trial slot is fully verified, then mark it confirmed, commit applicable
security counters, and finally allow the old slot to become reusable. A reset at
any point must yield either the old confirmed image or the new confirmed image,
never an unverified slot.

If health fails, the agent should record a compact durable diagnostic before
requesting rollback. The bootloader performs rollback even if the trial agent is
nonfunctional. After booting the recovery image, the agent reports
\code{ROLLED\_BACK} with the failed release ID and reason.

% -----------------------------------------------------------------------------
